\documentclass{article}
\usepackage[margin=1.5cm]{geometry}
\usepackage{amsmath}
\usepackage{enumitem}

\setlist{nosep,leftmargin=*}

\begin{document}
\twocolumn

\title{Chapter 11 Warm-Up: From Dictionaries to Classes}
\author{Prof. Jordan C. Hanson}
\date{}

\maketitle

\section{Review: Dictionary Basics}

A Python dictionary stores information as key--value pairs. For example,

\begin{verbatim}
students = {2601001: "Ana", 2601002: "Ben"}
\end{verbatim}
\noindent
Here, each student ID is a \textbf{key} and each name is its corresponding \textbf{value}.  What is printed by  \\ \texttt{print(students[2601002])}?  Add to the dictionary a student named Diego with ID 2601008.  A dictionary is useful when we want to associate one piece of information with another. A \textbf{class} goes further: it lets us define a new kind of object that can store several pieces of information and also define what that object can \textbf{do}.

\section{Complex Numbers}

A complex number has a \textbf{real part} and an \textbf{imaginary part}:

\begin{equation}
z = a + bj
\end{equation}

where

\begin{equation}
j^2 = -1
\end{equation}

To add complex numbers, add real parts together and imaginary parts together:

\[
(a+bi)+(c+di)
=
(a+c)+(b+d)i.
\]

Subtraction works the same way:

\[
(a+bi)-(c+di)
=
(a-c)+(b-d)i.
\]

\section{A Class for Complex Numbers}

We can create our own Python class that stores the real and
imaginary parts of a complex number.

\begin{verbatim}
class ComplexNumber:
    def __init__(self, real, imag):
        self.real = real
        self.imag = imag
\end{verbatim}

Now we can create two \textbf{instances}:

\begin{verbatim}
z1 = ComplexNumber(3, 2)
z2 = ComplexNumber(5, -4)
\end{verbatim}

For \texttt{z1}, the values are

\begin{verbatim}
z1.real  ->  3
z1.imag  ->  2
\end{verbatim}

In your own words, what does \texttt{self} represent?

\section{Teaching the Class to Add and Subtract}

Normally Python does not know how two objects of a new class should be added. We can define that behavior with the special method \texttt{\_\_add\_\_()}:

\begin{verbatim}
class ComplexNumber:
    def __init__(self, real, imag):
        self.real = real
        self.imag = imag

    def __add__(self, other):
        return ComplexNumber(
            self.real + other.real,
            self.imag + other.imag
        )
\end{verbatim}

Now Python understands

\begin{verbatim}
z3 = z1 + z2
\end{verbatim}

because \texttt{\_\_add\_\_()} tells Python what the \(+\) operator means for two \texttt{ComplexNumber} objects.  Now add a \texttt{\_\_sub\_\_()} method so that Python understands knows what the \(-\) means.

\section{Check Your Understanding}

\begin{enumerate}
\item Let $z_1 = 2+2j$, and $z_2 = 2-2j$.  Using your class, compute (a) $z_1 + z_2$, and (b) $z_1 - z_2$.
\item Let $z_1 = 3+0j$, and $z_2 = 0+3j$.  Using your class, compute (a) $z_1 + z_2$, and (b) $z_1 - z_2$.
\item How would you define a method within your \verb+ComplexNumbers+ class to handle \textit{multiplication}?
\end{enumerate}

\end{document}
